\section*{Bonding: communicate before committing}
Occupied carbon and oxygen show directional guides for valid neighbouring recipe sites. These are interface guides, not literal electron orbitals. A held compatible atom near a target produces a restrained electrical arc. A solid target ring and the words \textbf{Release to bond} confirm that releasing will attach it.

Initial tolerances are 18 cm for preview, 5 cm from the target for attachment, and 15 degrees from the authored bond direction. The direction is measured from the occupied atom toward the incoming atom, not from wrist orientation. These inspector settings are starting values for evaluation, not validated thresholds.

\begin{center}
\begin{tikzpicture}[node distance=13mm and 18mm]
\node[state] (held) {Held atom};
\node[state,right=of held] (near) {Compatible nearby\\arc preview};
\node[state,below=of near] (ready) {Position + angle valid\\release ring};
\node[state,below=of held] (free) {Invalid release\\still loose};
\node[state,below=of ready] (bond) {Release commits\\snap + sound};
\draw[flow] (held)--(near);\draw[flow] (near)--(ready);\draw[flow] (ready)--(bond);
\draw[flow] (held)--(free);
\draw[flow] (near.north)--++(0,.5)-| node[pos=.25,above,font=\scriptsize]{Leave preview area} (held.north);
\draw[flow] (ready.west)--node[above,font=\tiny]{Invalid release} (free.east);
\end{tikzpicture}
\par\small Figure 4. Bonding states; cancellation clears transient cues.
\end{center}


A shared evaluator decides both the preview and release outcome. A still-valid highlighted target remains stable if several targets are nearby. Release rechecks element, adjacency, vacant site, angle and distance; invalid attempts leave inventory and structure intact. Valid placement commits once, then animates a brief visual snap. Cancellation, reset, mode changes and tracking interruption clear active feedback.

\textbf{Scientific interpretation.} Covalent bonds are not electrical sparks. The tutorial explicitly identifies arcs as teaching cues. On completion, the transient effect disappears and solid rods retain the authored bond-order representation. The simulator remains recipe-based.

\pending{Observe at least one deliberately wrong-element release and one wrong-angle release. Record whether the learner can explain why each attempt failed, and whether the feedback predicts the next successful attempt.}
